% 子文件：8.多元函数微分学.tex
% 主文档使用 \input{8.多元函数微分学} 引入，不含 document 环境

\section{多元函数的基本概念}
\subsection{预备知识}

\subsection{多元函数}

定义8.1\quad 设$D$是$Oxy$平面的点集，若变量$z$与$D$中的变量$x,y$之间有一个依赖关系，使得在，记为

\subsection{多元函数的极限与连续}

设集合

定义8.8\quad 设$u$是$D$的聚点。

多元函数极限的含义是：

务必注意，虽然多元函数的极限与一元函数的极限定义相似，但它复杂得多。一元函数在某点处极限存在的充要条件是左右极限存在且相等。

顺便指出：$(x,y)\to(x_0,y_0)$的过程中，$x$和$y$是作为点的坐标。

定义8.9\quad 设函数$f(P)$的定义域为，且，则说函数。

若

有界闭区域上的多元连续函数有如下重要性质：
\begin{enumerate}
    \item 在
    \item 在
\end{enumerate}

\section{偏导数与高阶偏导数}
\subsection{偏导数}

定义8.10\quad 设函数$z=f(x,y)$在点，称为。若极限
\[
\lim_{\Delta x\to0}\frac{f(x+\Delta x,y)-f(x,y)}{\Delta x}
\]
存在，则称此极限值为关于$x$的偏导数，记为。同样定义处关于$y$的偏导数为。

如果在区域$E$内每一点，记为。同样，$z=$记为。

\noindent 偏导函数

对一般。

由偏导数的定义知，

一元函数可导必连续，但对多元函数，偏导数都存在，函数未必有极限。

因为偏导数。

\subsection{高阶偏导数}

记函数，它们仍是$E$内$x,y$的函数。，其中称为混合二阶偏导数。

递推地可以定义各阶偏导数，二阶。

定理8.1\quad 如果在点。

\section{全微分}

设函数

定义8.11\quad 若函数。

定理8.2\quad 若函数$z=f(x,y)$在点$P(x,y)$处可微，则在点$P$处偏导数$\dfrac{\partial z}{\partial x}$及$\dfrac{\partial z}{\partial y}$都存在，且
\[
\frac{\partial z}{\partial x}=A,\quad \frac{\partial z}{\partial y}=B.
\]
因为，我们把。

定理8.3\quad 若在点。

\section{复合函数求导法}

定理8.4\quad 如果函数在点
\[
\frac{\partial z}{\partial x}=\frac{\partial z}{\partial u}\frac{\partial u}{\partial x}+
\]

由上式。

同样条件下，公式可以推广到一般复合函数上去，设，则$z$关于某个自变量$x_j$的偏导数等于$z$对每个与$x_j$有关的中间变量$u_i$的偏导数与这个中间变量$u_i$对$x_j$的偏导数之积的总和，即。该公式称为链导法则。

在，当中间变量都是同一个自变量的一元函数时，，称为$z$的全导数公式。

链导法则公式可借助，称矩阵为函数，记为，则链导法则公式。

\noindent 全微分形式不变性。

\noindent 四则运算的全微分法则：

\section{隐函数求导法}
\subsection{一个方程的情况}

定理8.5（隐函数存在定理）\quad 设函数$F(x,y,z)$在点$(x_0,y_0,z_0)$的某邻域内具有连续的偏导数，且，则方程在点的某邻域内确定唯一一个函数，满足，且有偏导数公式
\[
\frac{\partial z}{\partial x}=-\frac{F'_x(x,y,z)}{F'_z(x,y,z)},\quad
\frac{\partial z}{\partial y}=-\frac{F'_y(x,y,z)}{F'_z(x,y,z)}.
\]

\subsection{方程组的情况}

设由方程组
\[
\begin{cases}
F(x,y,u,v)=0,\\
G(x,y,u,v)=0
\end{cases}
\]
确定两个二元函数$u=u(x,y),\ v=v(x,y)$。将恒等式
\[
\begin{cases}
F\bigl(x,y,u(x,y),v(x,y)\bigr)=0,\\
G\bigl(x,y,u(x,y),v(x,y)\bigr)=0
\end{cases}
\]
两边关于$x$求偏导，由链导法得到。解这个以$\dfrac{\partial u}{\partial x},\dfrac{\partial v}{\partial x}$为未知量的代数方程组，当系数行列式不为零时，即
\[
\begin{vmatrix}
\frac{\partial F}{\partial u} & \frac{\partial F}{\partial v}\\[4pt]
\frac{\partial G}{\partial u} & \frac{\partial G}{\partial v}
\end{vmatrix}\neq 0,
\]
解得。上面出现的由函数的偏导数构成的行列式称为雅可比行列式，记。于是。如果。

\section{偏导数的几何应用}
\subsection{空间曲线的切线与法平面}

设曲线$l$以参数方程给出。设，所以向量是切线的方向向量，称为曲线$l$在。故曲线的切线方程为。过点，其方程为。

\subsection{曲面的切平面与法线}

设曲面的隐函数方程给出，设其方程为。点，于是有。两边在，由全导数公式得。此式表明：向量与，切平面方程为。过，其方程为。现在。

\subsection{二元函数全微分的几何意义}

二元函数在点处可微。

\section{多元函数的一阶泰勒公式与极值}
\subsection{多元函数的一阶泰勒公式}

设$z=f(x,y)$在点$X_0(x_0,y_0)$的某邻域$U(X_0)$内有二阶连续的偏导数，$X(x,y)$是$U(X_0)$内的点。记$\Delta X=(\Delta x,\Delta y)=(x-x_0,y-y_0)$。称函数$z=f(x,y)$的两个一阶偏导数构成的向量$(f'_x,f'_y)_X$为$f(x,y)$的梯度。四个二阶偏导数构成的方阵
\[
H(X)=
\begin{pmatrix}
f''_{xx} & f''_{xy}\\
f''_{yx} & f''_{yy}
\end{pmatrix}_X
\]
称为$f(x,y)$的黑塞矩阵。

定理8.6\quad 若二元函数，则对，存在，使得
\[
f(x,y)=
\]
其中
\[
R=\frac1{2!}\bigl(\Delta x\quad \Delta y\bigr)
\begin{pmatrix}
f''_{xx} & f''_{xy}\\
f''_{yx} & f''_{yy}
\end{pmatrix}_{X^*}
\begin{pmatrix}
\Delta x\\
\Delta y
\end{pmatrix}.
\]
称上式为$f(x,y)$的一阶泰勒公式。

\noindent 推论\quad 在定理8.6的条件下。

\subsection{多元函数的极值}

定义8.12\quad 设二元函数，且，则说函数。

定理8.7（极值的必要条件）\quad 设函数在$X_0$处取极值，且在该点处函数的偏导数都存在，则必有
\[
\bigl(f'_x,f'_y\bigr)_{X_0}=\boldsymbol{0}.
\]

定理8.8（极值的充分条件）\quad 设是函数的驻点。若$f(X)$在点$X_0$的黑塞矩阵正定（负定），则为极小值（极大值）。

\noindent 推论\quad 设，记$f''_{XX}(x_0,y_0)$。
\begin{enumerate}
    \item 若
    \item
    \item
\end{enumerate}

\subsection{条件极值、拉格朗日乘数法}

设函数$z=f(x,y)$在约束条件$\psi(x,y)=0$下，在点$(x_0,y_0)$处取极值。

\section{方向导数和梯度}
\subsection{方向导数}

在数量场。

定义8.13\quad 设点，比式的极限存在，则称方向导数，记为，即。由定义。

定理8.9\quad 设在点$P_0$处可微，则函数，且，其中是$l$的方向余弦。

\subsection{梯度}

定义8.14\quad 数量场（函数）$u(P)$在点$P$处的梯度是个向量，其方向为$u(P)$在点$P$的变化率最大的方向，其模恰好等于这个最大的变化率，记为$\mathrm{grad}\,u$。

在空间直角坐标系下，梯度的表达式为。

可以直接推出下列性质及运算法则。
\begin{enumerate}
    \item
    \item
    \item
    \item
    \item
\end{enumerate}